\documentclass[10pt,openany]{book}
\usepackage{robbook}
\SetROBAccent{ROBRed}
\SetROBVolume{COMPLETE FIELD MANUAL - ADVANCED}{BUILDING R.O.B.}
\SetROBChapterPrefix{Chapter}

\begin{document}
\ROBTitle{Building R.O.B. - Complete Field Manual}{BUILDING R.O.B.}{The complete maker's field manual: fabrication, circuits, firmware, Mac software, perception, control, and operations}{Advanced evidence draft with builder placeholders}{2025-rob-side-full-body.jpg}

\ROBFrontMatter{An honest build record}{Complete in scope, intentionally incomplete where the archive cannot prove the detail.}{%
This manual connects ROB's 95-page design presentation, private construction-photo archive, Arduino base firmware, motor references, and current Mac, iPhone, Watch, Vision Pro, lidar, RGB-D camera, serial, MLX/VLM, Gemini, stage-show, arm, autonomy, H.264, and secure-control source trees.

It is a first evidence-based edition, not a stamped mechanical design, production wiring schematic, safety certification, or permission to energize the machine. Use the status labels on every subsystem. Replace each Builder input needed box with measured, reviewed information as the book iterates.

\begin{center}
\includegraphics[height=0.27\textheight]{rob-used-future-frontispiece.png}\\[-1pt]
{\scriptsize\itshape Original illustrated portrait inspired by ROB; it is visual storytelling, not engineering evidence.}
\end{center}
}

\section*{Evidence-status legend}

\begin{tabularx}{\textwidth}{p{1.75in} X}
\Implemented & Confirmed in the inspected current source or directly visible hardware evidence.\\[5pt]
\Commanded & Software sends a value, but this review did not measure the physical outcome.\\[5pt]
\Experimental & Present behind a gate, dependent on external hardware, or still under active development.\\[5pt]
\Legacy & Retained compatibility or historical behavior, not the preferred new design.\\[5pt]
\Planned & Described in a slide, TODO, protocol, or interface without proof of completed execution.\\[5pt]
\Unverified & The builder must measure, reconcile, or supply the missing record.\\
\end{tabularx}

\begin{safetybox}[title={NON-NEGOTIABLE SAFETY BOUNDARY}]
ROB combines large batteries, high-current brushless motors, chains, treads, flippers, linear motion, long arms, machine-tool fabrication, networking, cameras, and software-controlled movement. Qualified adults must own the hazard analysis, power protection, emergency-stop architecture, workholding, guarding, lifting, test perimeter, and public operation. Never infer a safe circuit or fabrication dimension from a photograph. Never treat an ordinary IR sensor, software timeout, or cloud model as safety-rated protective equipment.
\end{safetybox}

\ROBArduinoStatus
\ROBAmberStatus
\ROBSourceRoot

\section*{Source snapshot and publication cautions}

The initial evidence pass began on 2 August 2026; the consolidated Cerebro, Vision Pro, local-model, stage-show, AMBER, and advanced video pass reflects the working trees inspected through 11 August 2026. Active changes still require an explicit tag or commit before publication. Original gallery photographs were converted with orientation applied and metadata stripped. Attendee photographs, especially images containing children, require release review or careful exclusion.

The archived B-COMMAND document prohibits reproduction without permission, and the HengDrive outline is marked proprietary and confidential. This manual does not reproduce or paraphrase their specifications. It mentions the motor-family label visible in a builder-owned photograph only to identify the open evidence gap. Obtain permission and a current authorized public datasheet before publication.

\tableofcontents
\clearpage
\mainmatter

\part{Intent and architecture}

\ROBChapterPhoto{Chapter 1}{The robot as a long-running experiment}{2025-rob-face-portrait.jpg}

\chapter{Design history and current scope}

ROB is not one frozen machine. The archive shows a succession of working prototypes, public demonstrations, repair decisions, new sensors, drivetrain modifications, camera changes, arm work, and software migrations. A useful build manual must preserve that timeline because the answer to ``How is it built?'' depends on the date and subsystem.

\section{Evidence-based chronology}

\begin{tabularx}{\textwidth}{>{\robcondensed\bfseries\color{ROBAccent}}p{0.72in} X X}
\toprule
PERIOD & VISIBLE OR DOCUMENTED STATE & EDITORIAL INTERPRETATION\\
\midrule
2019 & Functional black tracked robot with screen torso, earlier camera head, exposed arms, and movement demonstrations. & An integrated early prototype existed; exact internal revision and specifications are not documented.\\
2020-2021 & Open electronics, camera/head revisions, wooden base templates, tread geometry, battery layout, and plans for CNC aluminum. & The builder was using accessible materials to resolve packaging and mechanics before a final enclosure.\\
2022 & Mini computer and serial cabling, base brackets, drilling, actuator hardware, arm electronics, OAK-D camera, and public appearances. & Major integration work joined mechanics with onboard computing and sensing.\\
2023 & Public use and head/camera refinements, with fewer construction photographs. & Operational experience likely drove requirements, but the archive does not link each change to a written test.\\
2024 & Sprocket-area machining, chain clearance, tread measurement, internal rewiring, batteries, actuator, voltage gauge, and controller board. & A substantial drivetrain and base-serviceability revision is strongly documented.\\
2025 & Shoulder and neck gearing, arm removal/reassembly, cable labels, emergency-stop hardware, Mac software, vision, mapping, and chess reach studies. & The current visual identity and advanced software experiments are visible, while final calibration remains in progress.\\
\bottomrule
\end{tabularx}

\ROBPhotoTwo{2019-rob-front-corner.jpg}{The 2019 machine already joins a tracked base, arms, sensors, and operator-facing character.}{2025-rob-side-full-body.jpg}{A 2025 side view shows the newer spherical camera head, speaker face, neck mechanism, articulated arm, and tracked base.}

\section{Design principles that recur}

Across the sources, several durable principles emerge:

\begin{itemize}
  \item use prototypes to discover packaging and load problems;
  \item separate high-level intent from low-level actuator values;
  \item keep independent stop layers close to each source of failure;
  \item isolate optional cameras, Python helpers, video, speech, and AI from core motion control;
  \item preserve service access and label wiring;
  \item make unsupported capability fail honestly;
  \item collect evidence before calling a proposal implemented.
\end{itemize}

\section{Current subsystem status at a glance}

\begin{tabularx}{\textwidth}{p{1.42in} p{1.4in} X}
\toprule
SUBSYSTEM & STATUS & EVIDENCE-BASED SUMMARY\\
\midrule
Tracked base & \Implemented & Physical treads, chains, drive motors, brackets, batteries, and actuator are visible across multiple revisions. Exact final dimensions and performance are missing.\\
Arduino base & \Implemented & Historical firmware parses a seven-field frame and commands treads, flipper, and actuator; it reads IMU and IR data. It has significant safety and parser defects.\\
Flipper motion & \Commanded & Firmware and controller paths exist. Final geometry, limits, load rating, and measured motion are unverified.\\
Torso / third-wheel actuator & \Commanded & Serial control and physical actuator evidence exist, but the presentation assigns conflicting roles across revisions.\\
Mac mini / Cerebro & \Implemented & Central macOS coordination, serial, camera, speech, secure control, and autonomy infrastructure exist in source.\\
Phone controller & \Implemented & Established control UI, pairing, autonomy, and action-console paths exist.\\
Watch controller & \Implemented & Freshness-checked motion and speech intent relay through the phone.\\
Vision Pro controller & \Experimental & Secure control adapter, simulator, leases, and separate H.264 viewing path exist; hardware deployment still needs release validation.\\
RGB-D camera & \Experimental & Supervised DepthAI helper and aligned RGB/depth transport exist; device and environment validation remain necessary.\\
RPLidar & \Experimental & Typed scan/map publication exists; native Mac deployment of the proprietary framework remains unresolved.\\
Arms / neck & \Commanded & Servo, Amber script, and keyframe paths exist; calibration, feedback, limits, and collision safety are incomplete.\\
Gemini actions & \Experimental & The multimodal session defaults off, while audio/video default on once that session is enabled unless separately disabled. The action tool is separately default-off; only \texttt{stop\_motion} has a partial local executor.\\
\bottomrule
\end{tabularx}

\ROBPlaceholder{Builder narrative: add the project's origin, meaning of R.O.B., design goals, major failures, people and shops who contributed, dates of each hardware revision, and which public demonstrations changed the design.}

\ROBChapterPhoto{Chapter 2}{System boundaries and ownership}{2024-open-tracked-base.jpg}

\chapter{Whole-robot architecture}

\ROBControlStackDiagram

\section{Physical and information paths}

The robot contains two overlapping architectures:

\begin{enumerate}
  \item The \textbf{energy architecture} moves stored electrical energy through protection, conversion, controllers, conductors, and motors.
  \item The \textbf{information architecture} moves intent and measurements through interfaces, authentication, state machines, protocols, adapters, firmware, and feedback.
\end{enumerate}

The two paths meet at motor controllers and actuators. Keeping their documentation separate prevents a common error: assuming that a correct data value proves the high-current circuit is safe, or that a correctly wired motor proves the command path is authorized.

\section{Desired component ownership map}

The table below states the target responsibility split. It is a design contract, not a claim that every current implementation already satisfies every responsibility; gaps must remain visible in the subsystem status and verification chapters.

\begin{tabularx}{\textwidth}{>{\robcondensed\bfseries}p{1.32in} X X}
\toprule
OWNER & PRIMARY RESPONSIBILITY & MUST NOT ASSUME\\
\midrule
Human operator & Select mode, maintain line of sight and perimeter, respond to hazard, own start and stop. & That autonomy, sensors, or remote viewers replace supervision.\\
Controller app & Capture fresh intent, apply its client-specific lease or watchdog, present authoritative mode and stop state. & That sending a packet proves movement occurred.\\
Secure transport & Authenticate device, authorize role, bound and order messages, reject replay and stale sessions. & That an authenticated controller's payload is physically safe.\\
Cerebro & Arbitrate control master, coordinate subsystems, validate high-level actions, adapt to hardware protocols. & That optional camera, AI, or lidar availability can weaken manual stop.\\
Arduino base & Target: accept only complete bounded base commands, own a local deadline, and drive low-level pins and actuator serial. & That the current historical parser already meets this target, or that any arriving byte is a valid command or heartbeat.\\
Motor electronics & Enforce rated electrical behavior, current limits, and safe-start rules. & That software alone prevents wiring or load faults.\\
Mechanical system & Carry loads, constrain motion, guard hazards, survive declared environments. & That commanded direction, speed, or stop equals measured behavior.\\
\bottomrule
\end{tabularx}

\section{Sense-plan-act without surrendering authority}

\ROBSimpleSystemDiagram

Manual control, deterministic autonomy, and AI-assisted behavior can all use the same physical plant, but they should enter through explicit authority gates. \Experimental\quad Current social-roam autonomy consumes fresh filtered lidar scans and, near 82 percent of its configured activation radius, attempts to turn and drive toward its origin. That behavior is not a physical geofence and does not prove containment. Optional Gemini tools declare high-level action proposals, but most do not yet have executors. Neither path should bypass local validation, the control-master state, the Arduino deadline, or the physical emergency stop.

\begin{advancedbox}
Document safety as a control hierarchy, not a list of features. For every hazard, name the prevention, detection, mitigation, independent stop, diagnostic coverage, proof test, and residual risk. Ordinary cameras, IR sensors, and software timers may support risk reduction but are not automatically safety-rated.
\end{advancedbox}

\ROBPlaceholder{Create an author-approved current block diagram showing installed batteries, disconnect and protection, DC buses, converters, Mac mini, any Raspberry Pi units, the Arduino-compatible base controller with exact model, motor controllers, Maestro, Amber arm controllers, cameras, lidar host, router, and every external port. Mark planned and removed hardware separately.}

\part{Mechanical build}

\ROBChapterPhoto{Chapter 3}{Chassis, treads, and drivetrain}{2024-nearly-assembled-base.jpg}

\chapter{Build the mobile base}

The base is the foundation for every later system. Presentation pages 3-9 show left and right wooden side templates, motor openings, wheel mounts, a battery opening, and the desire to reproduce the form in black powder-coated aluminum. Pages 10-20 document the tread tensioner, wheels, sprockets, and shaft collars. Later gallery photographs show extensive metal and wiring revisions.

\section{Start from interfaces}

Before drawing the final plate, freeze or explicitly version these interfaces:

\begin{itemize}
  \item tread-frame mounting points and full motion envelope;
  \item drive and idler shaft positions, diameters, retention, and bearing supports;
  \item motor face, shaft, and cable clearances;
  \item chain line, sprocket sizes, guard volume, and service access;
  \item flipper shaft, torque path, rear/front location, and end stops;
  \item battery envelope, restraint, protection access, and removal path;
  \item actuator mounts and full extension/retraction envelope;
  \item lifting, tie-down, stand, cover, and inspection points.
\end{itemize}

\ROBPhotoTwo{2022-drive-assembly-rear.jpg}{A centered rear view shows the drive assembly between two treads.}{2024-sprocket-machined-opening.jpg}{A later sprocket view shows why chain-line and plate clearance need dimensioned control.}

\section{Tread tension and support}

The presentation describes aluminum bars with threaded holes that pull the tread assembly outward, free-spinning wheels retained on shafts, and a heavy-load problem that led to improvised HDPE and skateboard-bearing supports. These are valuable design clues but not a final service specification.

Record a repeatable tension method. A statement such as ``tight enough'' cannot be audited. Depending on the tread system, the specification might be deflection under a known force, tensioner position, exposed thread length, torque plus verification, or the manufacturer's procedure.

\ROBPhotoTwo{2022-chain-sprocket-interior.jpg}{A large sprocket and chain are visible inside the base.}{2024-chain-drive-closeup.jpg}{The close chain path illustrates alignment, abrasion, and guarding concerns.}

\begin{safetybox}
Chains, sprockets, and treads create drawing-in, nip, shear, and entanglement hazards. Disconnect and verify power, block stored motion, prevent unintended rolling, and use fixed guards before hands enter the mechanism.
\end{safetybox}

\section{Flipper revision and third-wheel concept}

Presentation pages 21-26 say the flippers moved from the front wheels to the rear after a design revision. The deck also records a one-degree bar offset, cotter-pin concerns, and a proposed perpendicular third wheel raised by an actuator to reduce tread scrub while turning.

\ROBPhoto[0.54\textheight]{rob-v3-page-23.jpg}{The annotated upside-down prototype on presentation page 23 captures the third-wheel idea and identifies a design revision. It is not a final drawing.}

\Planned\quad Treat the third wheel as a separate design project until the builder supplies the geometry, load path, stability analysis, actuator selection, ground clearance, deployed/retracted sensing, failure response, and test evidence.

\section{Provisional assembly sequence to verify}

\begin{enumerate}
  \item inspect plate dimensions, datum features, threads, deburring, and finish;
  \item install bearings, shaft supports, and fixed structural brackets;
  \item trial-fit tread frames, drive motors, shafts, sprockets, and chains without power;
  \item measure shaft parallelism, sprocket coplanarity, tension range, and full-motion clearance;
  \item install flipper shaft and drivetrain with temporary manual limits;
  \item fit actuator brackets through the full stroke with the load supported independently;
  \item install guarded power components and labeled harnesses only after mechanical inspection;
  \item add batteries using a rated lift and restraint plan;
  \item perform continuity, insulation, polarity, and protective-device checks;
  \item test one actuator at a time on stands under a bounded current and command limit;
  \item install guards and covers before floor testing.
\end{enumerate}

\ROBPlaceholder{Supply final chassis length, width, plate thicknesses and alloys, tread model, wheel and sprocket data, shafts, bearings, collars, chain, tension procedure, flipper dimensions, fastener schedule, torque values, guard drawings, total mass, center of mass, and as-built CAD.}

\ROBChapterPhoto{Chapter 4}{Torso, body, head, and arms}{2025-rob-side-full-body.jpg}

\chapter{Build above the base}

The upper structure must hold cameras and arms while surviving acceleration, arm moments, cable pull, transport, and public interaction. The presentation openly records weak prototype areas: body-to-torso joints pulling on metal, bent L brackets, and wooden parts glued and screwed together.

\section{Body-to-torso connection}

Presentation pages 56-60 show the connection under load and call for mechanical improvement. A final design should provide:

\begin{itemize}
  \item a defined load path into the chassis;
  \item bearing surfaces or pivots sized for radial, axial, and moment loads;
  \item hard stops and motion limits independent of software;
  \item retention that remains safe after one fastener loosens, where practicable;
  \item inspection access and visible wear indicators;
  \item cable routing that does not become the motion stop;
  \item a transport lock that does not depend on powered holding torque.
\end{itemize}

\section{Body rotation and cable continuity}

Presentation page 67 proposes a 36 mm motor and a body bearing with rotation constrained from $-\pi$ to $+\pi$ instead of using a slip ring. A nearby restricted B-COMMAND file exists, but its contents are excluded and the firmware and photographs do not prove that the associated component is installed.

\Planned\quad A rotating body requires a controlled cable-wrap budget or an approved rotary interface, home/index sensing, hard and soft limits, collision envelopes, and a recovery plan after position loss.

\ROBPhotoTwo{2025-rear-electronics-bay.jpg}{The rear bay shows how upper-body motion can pull on a dense harness.}{2025-labeled-cables.jpg}{Temporary labels aid teardown, but permanent routing must account for repeated flex and abrasion.}

\section{Head and neck}

The archive spans an earlier spherical multicamera head, an Insta360 Pro II concept, OAK-D sensors, a later spherical head, and a 2025 geared neck mechanism. The presentation calls the earlier hobby-servo arrangement inadequate and asks for tilt plus full rotation.

\ROBPhotoTwo{2022-oakd-depth-camera.jpg}{The OAK-D stereo camera is a distinct sensing module, not the whole head design.}{2025-neck-gear-train.jpg}{The 2025 neck gear train exposes stacked gears and support plates that need calibrated limits and guarded pinch points.}

\section{Arms}

The deck proposes two Amber seven-degree-of-freedom arms with force/current-sensing grippers. Current software includes Amber Python/SSH paths, Maestro servo control, and keyframe fields. Photographs show the installed articulated arms and their controller boards. This establishes integration work, not a verified payload or autonomous grasp capability.

\ROBPhotoTwo{2025-detached-articulated-arm.jpg}{A detached arm reveals its many joints and long reach.}{2022-arm-controller-board.jpg}{A circular controller board sits behind the arm mechanism.}

\Commanded\quad Arm and neck values can be issued through current control paths.\par
\Unverified\quad Joint feedback, zeroing, limits, load rating, collision envelope, emergency behavior, and physical outcome were not measured in this review.\par
\Planned\quad General object pickup remains unavailable because camera-to-arm transforms, inverse kinematics, collision checking, and a verified grasp executor are incomplete.

\ROBPlaceholder{Supply upper-body drawings; body bearing and motor; torso/third-wheel actuator roles; head sensor inventory; neck axes, ratios, limits, and home sensors; arm model/firmware, joint definitions, zeroing, payload, reach, current limits, gripper behavior, collision envelope, cable routing, and transport restraints.}

\ROBChapterPhoto{Chapter 5}{From template to finished part}{2022-drilling-metal-rail.jpg}

\chapter{Fabrication and finish}

The presentation asks how to print and weld the aluminum, whether templates or a 3D scan would help, and whether an outside robotics shop should quote the metalwork. Those questions prove that a complete final fabrication recipe is not present. This chapter supplies the record structure the builder should fill.

\section{Measurement plan}

Choose datums before collecting dimensions. For every critical measurement, record:

\begin{itemize}
  \item feature and datum names;
  \item nominal value and unit;
  \item tolerance or functional clearance;
  \item instrument and resolution;
  \item repeated measurements and uncertainty;
  \item temperature or load state if relevant;
  \item photograph or drawing revision;
  \item who measured and when.
\end{itemize}

\ROBPhotoTwo{2024-caliper-tread-measurement.jpg}{A digital caliper is used across the tread frame; the publication needs the datum and recorded value.}{2022-bracket-fit-check.jpg}{A fit check can reveal interface errors before a hole becomes permanent.}

\section{Process traveler}

For each fabricated part, create a traveler with these gates:

\begin{enumerate}
  \item released drawing and material certificate or stock identification;
  \item saw/cut blank with allowance;
  \item workholding setup and datum pickup;
  \item roughing operation;
  \item intermediate inspection;
  \item finish holes, pockets, threads, bends, or weld preparation;
  \item deburr and edge-break requirement;
  \item weld/braze/bond sequence and distortion check, if applicable;
  \item surface preparation and coating;
  \item final dimensions, photographs, and acceptance signature;
  \item nonconformance and rework record.
\end{enumerate}

\begin{safetybox}
The builder or qualified shop must write the actual machine, cutter, feeds and speeds, workholding, coolant, chip control, guarding, ventilation, hot-work, welding, grinding, coating, and PPE instructions. A generic book cannot make a particular shop operation safe.
\end{safetybox}

\section{Design for assembly and service}

Use captive or accessible fasteners where covers need regular removal. Separate abrasive motion and sharp edges from harnesses. Provide inspection windows for chain line, tension, battery restraints, and protective devices. Ensure a heavy component can be supported before its final fastener is removed. Include lifting points and a stand plan in the CAD, not as an afterthought.

\ROBPlaceholder{Insert the builder-approved machining and assembly travelers, final CAD/template files, CNC or manual process, weld procedure, coating system, adhesives, inspection tools, rejected-part lessons, outside vendors, costs, and photographs paired with each actual operation.}

\part{Power, electronics, and firmware}

\ROBChapterPhoto{Chapter 6}{Power architecture}{2024-batteries-and-power-leads.jpg}

\chapter{Store, convert, protect, and disconnect energy}

\section{What the archive says}

The 2021-2022 presentation describes four 20 Ah LiFePO4 batteries for a stated total of 80 Ah and shows a 12 V distribution bus, an onboard charger concept, a 120 VAC inverter, and a 12-to-24 V 240 W converter. Later photos show a different visible battery arrangement, including two large packs. The current topology is not proven.

\ROBPhotoTwo{rob-v3-page-41.jpg}{Historical 2021--2022 presentation concept: page 41 sketches four batteries and packaging around the treads; it is not evidence of the current power topology.}{2024-batteries-and-power-leads.jpg}{The 2024 photograph shows the physical packs and leads visible in that revision.}

\section{Required single-line diagram}

The final manual needs an electrical single-line diagram from every energy source to every load. At minimum, identify:

\begin{itemize}
  \item battery pack chemistry, nominal and full-charge voltage, capacity, BMS, and series/parallel connection;
  \item main fuse or breaker near each energy source;
  \item main disconnect and emergency-stop contactor/controller path;
  \item precharge, inrush, anti-spark, or safe connector sequence where required;
  \item 12 V, 24 V, 5 V, and other buses with converter ratings and derating;
  \item tread, flipper, actuator, arm, compute, sensor, light, audio, router, and accessory branches;
  \item wire gauge, insulation rating, connector, polarity, routing, and strain relief;
  \item grounding, chassis bonding, shielding, and separation of noisy and sensitive returns;
  \item charger interlock, shore-power protection, inverter isolation, and exposed receptacles;
  \item test points, indicators, and safe discharge behavior.
\end{itemize}

\ROBPhotoTwo{2024-power-cable-pass-through.jpg}{A pass-through needs abrasion protection and restraint.}{2024-voltage-gauge.jpg}{A 13.1 V display is a diagnostic clue, not proof of state of charge or system health.}

\section{Motor electrical evidence}

The motor label visible in build photographs identifies a HengDrive B5685G family, but the archived vendor sheets carry restrictive or confidential notices. This public draft therefore does not reproduce or paraphrase their specifications. Obtain written permission or current public datasheets, then identify each installed winding, gearbox, controller, wire function, current limit, branch protection, and measured operating envelope.

\begin{advancedbox}
A presentation slide proposes a 240 W converter. At 24 V, 240 W corresponds to 10 A before losses, but the installed topology and authorized motor ratings are not established. This does not approve or reject the circuit; it demonstrates why a measured current and power budget is mandatory.
\end{advancedbox}

\section{Emergency stop and safe energy state}

A physical E-stop should be independent of the UI and ordinary network. Its design must follow the actual hazard analysis: which energy is removed, which holding brakes remain, whether abrupt removal creates a worse hazard, how reset works, how contact welding or controller failure is detected, and how the operator proves it before a show.

\ROBPhoto{2025-emergency-stop-bench-layout.jpg}{A yellow emergency-stop button and control boxes appear in a 2025 bench layout. The photograph does not reveal contacts, safety category, or final wiring.}

\ROBPlaceholder{Provide the authoritative as-built schematic and calculations: present battery packs and topology, BMS, chargers, fuses/breakers, disconnect, E-stop circuit and contact ratings, converters, branch currents, wire gauges, connectors, grounding/bonding, inverter/shore-power safety, runtime, thermal measurements, and fault tests.}

\ROBChapterPhoto{Chapter 7}{Arduino base controller}{2024-dense-internal-routing.jpg}

\chapter{Base electronics and firmware}
\ROBSourceNow{ROBArduino/ROBOT_CEREBELLULAR_BASE_APP/ROBOT_CEREBELLULAR_BASE_APP.ino}{This is the current firmware reference for the entire chapter. Torso and Head sketches are cited only where their retired behavior is being compared.}

\section{Three historical roles, one present role}

The \texttt{ROBArduino} archive resolves an important architecture question. Earlier ROB revisions had three named Arduino sketches, but the builder reports that present-day ROB uses only the Base code.

\begin{tabularx}{\textwidth}{>{\robcondensed\bfseries}p{0.72in} p{1.45in} X}
\toprule
ROLE & PUBLICATION STATUS & INSPECTED SOURCE BEHAVIOR\\
\midrule
Base & \Implemented\ current role reported & Seven-field tread/flipper/actuator command; IR and MPU9250 processing; USB host link and actuator SoftwareSerial. The exact flashed binary and board remain to be recorded.\\
Torso & \Legacy\ retired & Fourteen four-character Maestro targets for channels 0, 1, and 6--17; MPU9250; disabled GPS path. Its timeout actuation is commented out. Two archived filenames contain identical bytes.\\
Head & \Legacy\ retired & MPU9250 and PIR experiment with disabled GPS declarations. The active parser consumes only one copied legacy field, not the Base seven-field command; timeout stop actuation is commented out.\\
\bottomrule
\end{tabularx}

Current diagrams, startup procedures, and hardware counts must show one Base Arduino. Historical diagrams may show all three only when labeled with their revision context. Cerebro still contains serial names and helpers for Base, Torso, Head, and Maestro; those software paths are evidence of compatibility history, not proof that three Arduino boards remain installed.

\section{Confirmed pin and communication map}

The firmware targets an Arduino-compatible controller, and historical controller logs call it an Arduino Mega. The exact board model and revision, core, and library versions remain to be supplied.

\begin{tabularx}{\textwidth}{>{\robcondensed\bfseries}p{1.38in} >{\centering\arraybackslash}p{0.7in} >{\centering\arraybackslash}p{0.7in} >{\centering\arraybackslash}p{0.78in} X}
\toprule
LOAD & SPEED & BRAKE & DIRECTION & NOTES\\
\midrule
Right tread & D2 & D3 & D25 & Active-low speed mapping in firmware\\
Left tread & D4 & D5 & D27 & Mirrored direction polarity\\
Flipper & D6 & D7 & D29 & No implemented limit input\\
\bottomrule
\end{tabularx}

\begin{tabularx}{\textwidth}{>{\robcondensed\bfseries}p{1.55in} p{1.2in} X}
\toprule
FUNCTION & CONNECTION & CURRENT SOURCE BEHAVIOR\\
\midrule
Mac base serial & USB, 250,000 baud & Fixed-width command plus diagnostic text and an existing role-specific startup line\\
Actuator controller & SoftwareSerial D22/D23, 19,200 baud & Compact serial forward/reverse; nominal speed -3200 to +3200; special inputs +3201 and -3201 clear safe start before commanding the corresponding 3200 magnitude\\
IR array & A5 through A0 & Six SharpIR instances; front/back warning thresholds at 25 cm; numeric six-sensor printout is commented out\\
IMU & I2C via Wire; D20/D21 if Mega & MPU9250 initialization and orientation telemetry\\
Speed feedback & Proposed D18/D19/D20 & Commented TODOs disagree: the header calls D18 left and D19 right, while later pin declarations reverse them; D20 conflicts with Mega I2C SDA\\
\bottomrule
\end{tabularx}

Cerebro consumes the existing \texttt{WARNING! FRONT}, \texttt{WARNING! BACK}, and blocking-error sentences without requiring a firmware flash. Its SceneKit diagnostic highlights the corresponding pair for three seconds. Silence becomes ``clearance unknown,'' not ``path clear,'' because the legacy stream has neither explicit clear events nor numeric distances. These terrain-sensitive IR warnings do not replace RPLidar traversability or a physical stop.

The existing Base sketch emits \texttt{BEGIN BASE STARTUP SEQUENCE} after IMU initialization. Cerebro pulses DTR to restart each USB serial candidate, waits up to 15 seconds for that exact role-specific line, retains the matching descriptor, and reruns the scan when the operator refreshes. It transmits no command bytes and requires no firmware flash. This removes reliance on a hub-dependent \texttt{/dev/cu.*} name, but it is role discovery rather than cryptographic authentication or hardware verification.

\section{Serial frame}

\begin{center}
\begin{tcolorbox}[colback=ROBInk,colframe=ROBAccent,width=0.96\textwidth,arc=0mm]
{\ttfamily\color{white}\large
\textcolor{ROBAccent}{\textasciitilde}+0000,+0100,+0000,+0100,+0000,+0000,+0000}
\end{tcolorbox}
\end{center}

The 42 ASCII bytes represent left brake, left signed speed, right brake, right signed speed, flipper brake, flipper signed speed, and actuator signed speed. Brake field 0 maps to HIGH and is described as released; nonzero maps to LOW and is described as engaged. Tread and flipper speed output is \texttt{255 - abs(command)}, consistent with the active-low motor control concept. The physical brake electronics are not documented by the five-wire motor outline.

\section{Actuator compact serial}

The code sends 0x83 to exit safe start, 0x85 for forward, and 0x86 for reverse. It encodes magnitude in a five-bit low chunk and remaining high bits. Special input values +3201 and -3201 also clear safe start before becoming +3200 or -3200. Reset should be separated from speed in a future typed protocol.

\section{IR and IMU status}

\Experimental\quad Six IR sensors are constructed and front/back readings are sampled. After ten repeated close readings, the code prints a warning. The two lines that assert movement-inhibit flags are commented out, side sensors do not participate, and full six-value telemetry is commented. The builder reports grass/floor false triggers.\par

\Implemented\quad MPU9250 acceleration, gyro, and magnetometer readings feed a Mahony quaternion filter and print yaw, pitch, roll, and temperature. The orientation is telemetry only; tilt does not inhibit motion. The yaw correction is hard-coded for an older Austin, Texas declination and must not be treated as current calibration.

\section{Safety and correctness audit}

The historical firmware should be preserved as an artifact, then replaced behind fixtures. Priority defects include:

\begin{enumerate}
  \item the five-byte numeric buffer lacks a C-string null terminator before \texttt{atoi} and debug printing;
  \item one available byte triggers 41 further reads without proof of a complete frame;
  \item separator positions, signs, ranges, checksum, and frame completeness are not validated;
  \item an extra read after the seventh field can consume the next frame's start marker;
  \item any received byte refreshes the deadman before frame validation;
  \item the deadman depends on loop counts rather than elapsed monotonic time;
  \item timeout sends zero speed with brake fields that release brakes;
  \item IR inhibit assignments are disabled and would not update outputs until another command path runs;
  \item one rotation test duplicates the same condition and omits the mirrored turn;
  \item command values are not clamped to PWM and actuator ranges;
  \item ramp constants, encoder feedback, and voltage monitoring are TODO-only;
  \item brake, direction, and PWM update ordering has not been proven jerk-free;
  \item the flipper and actuator have no firmware-visible end stops, position feedback, or motion timeout;
  \item setup writes the inactive value 255 to both tread PWM pins but never initializes flipper PWM D6; after \texttt{pinMode(OUTPUT)}, its default LOW is the maximum-command level under the active-low mapping, so boot safety relies on the unverified flipper-brake path;
  \item IMU initialization can abort before motor outputs are explicitly initialized safe;
  \item direction method names, comments, and debug strings disagree.
\end{enumerate}

\begin{safetybox}
Do not ``clean up'' direction logic on the live robot by inspection. First capture protocol fixtures, lift and guard the drivetrain, limit current and speed, verify every brake and direction state, measure timeout behavior, and retain an immediate independent stop.
\end{safetybox}

\section{Replacement-firmware architecture}

A modern base controller should contain:

\begin{itemize}
  \item safe GPIO initialization before optional sensors;
  \item bounded ring buffer and explicit framed decoder;
  \item version, length, type, sequence, range, and checksum validation;
  \item atomic command acceptance only after a complete valid frame;
  \item a monotonic deadline refreshed only by a valid authorized command;
  \item explicit states: boot-safe, disarmed, armed-neutral, moving, stopping, fault, E-stop;
  \item separate commands for reset, arm, drive, flipper, actuator, and diagnostic requests;
  \item clamped outputs, ramp and direction-change interlocks, and defined brake/coast policy;
  \item end-stop, encoder, current, voltage, and temperature feedback where hardware supports it;
  \item structured telemetry with units and status flags;
  \item unit, parser-fuzz, timing, and hardware-in-the-loop tests.
\end{itemize}

\ROBPlaceholder{Confirm board model and core; installed libraries; every connector and logic voltage; brake circuit; motor/flipper polarity; allowed command ranges; measured deadman deadline; brake/coast behavior; encoder status; IR models/positions/calibration; IMU orientation; actuator controller, end stops, current limits, and safe state.}

\part{Software and perception}

\ROBChapterPhoto{Chapter 8}{Cerebro on macOS}{2022-fanless-computer-and-serial.jpg}

\chapter{Central coordination on the Mac mini}
\ROBSourceNow{Cerebro/Cerebro/ROBMainViewController.mm}{This chapter begins at Cerebro's coordinator, then names each helper source as ownership crosses into serial, camera, autonomy, AI, or animation code.}

Cerebro is a mixed Swift, Objective-C, and Objective-C++ macOS application. Its central view/controller layer coordinates UI state while specialized components own serial, camera, speech, autonomy, security, and animation.

\section{Principal components}

\begin{tabularx}{\textwidth}{>{\ttfamily\scriptsize}p{2.35in} X}
\toprule
SOURCE COMPONENT & RESPONSIBILITY\\
\midrule
ROBMainViewController.mm & Operator UI, controller messages, action/autonomy surfaces, subsystem coordination.\\
ROBSerialBox.m / .h & Present Base serial ownership and heartbeat rendering, plus legacy discovery/helpers named for retired Torso and Head Arduino roles and Maestro hardware.\\
CameraManager.swift & Camera providers, frame-set ownership, bounded delivery, RGB-D helper lifecycle.\\
CameraViewController.swift & Vision processing and camera presentation.\\
ROBSpeechBox.m & Local speech recognition configuration and speech synthesis.\\
ROBAutonomyCoordinator.swift & Bounded session state and deterministic local autonomy coordination.\\
ROBAI.swift & Optional multimodal Gemini session and high-level tool coordination.\\
GeminiRoboticsProtocol.swift & Versioned messages and validation for model-facing robot actions.\\
KeyframeAnimationManager.swift & JSON-persisted animation values for arms, head/neck, treads, flipper, actuator, torso, and speech.\\
AutoNet shared/transport files & ROBControl v2 framing, pairing, roles, session identity, and authorization.\\
\bottomrule
\end{tabularx}

\section{Base heartbeat ownership}

Cerebro evaluates controller snapshots on a timer running at about 10 Hz and treats a snapshot as stale once its local receipt age exceeds 0.6 seconds. On the first transition from fresh to stale, it writes one neutral/braked serial frame, stops USB refresh, releases control authority, and allows the Arduino heartbeat to expire. This is a time-based gate, not a fixed count of three missed messages. The layered behavior is stronger than either deadline alone, but the Arduino's own elapsed time and physical stop result still require measurement.

\section{ROBSerialBox as an integration history}
\ROBSourceNow{Cerebro/Cerebro/ROBSerialBox.m}{Read \texttt{initialize\_connection}, \texttt{usbSerialPortPaths}, \texttt{connectToDetectedBase}, \texttt{probeBaseFirmwareAtPath}, \texttt{consumeBaseSerialBytes}, the Vision neck/torso methods, and \texttt{stopBaseMotionAndDropHeartbeat} together.}

\texttt{ROBSerialBox} is the current physical-I/O façade and an archaeological record. It contains Base, Maestro, Tic, AMBER, controller, and retired Head/Torso surfaces. At startup every POSIX descriptor receives the closed sentinel $-1$. Base discovery enumerates only USB callout devices, opens candidates at 250,000 baud, pulses DTR, sends no probe command bytes, and passively waits up to 15 seconds for the already-flashed firmware line \texttt{BEGIN BASE STARTUP SEQUENCE}. Hub-dependent \texttt{/dev/cu.*} names therefore need not be stored. The banner identifies a firmware role; it does not authenticate the board or verify wiring.

Serial reads are arbitrary byte chunks, not messages. Cerebro accumulates bounded data, extracts newline-delimited records, and then interprets them. It recognizes the existing front/back IR warning sentences and a bounded six-value numeric form when available. Legacy warning silence is reported as clearance unknown. Maestro selection instead uses USB vendor/product/interface identity, while the waist Tic path uses preflighted \texttt{ticcmd}; device discovery is protocol-specific rather than ``first port wins.''

The next refactor should retain the existing Objective-C façade for Interface Builder while moving each descriptor, read loop, write queue, reconnect generation, and parser into a single-owner component. Checking \texttt{fd != -1} cannot by itself make a concurrent close/write atomic. This is the practical lesson behind historical \texttt{Bad file descriptor} floods.

\section{External helpers and dependencies}

The OAK-D path uses a supervised Python helper rather than loading DepthAI directly into the control process. Cerebro can select or create a Python environment and validate the required import. Optional Amber scripts and tools such as \texttt{sshpass} and \texttt{ticcmd} are preflighted; missing dependencies disable their feature instead of terminating the app.

\begin{advancedbox}
Subprocess boundaries create explicit failure and update surfaces. The parent must validate executable paths, avoid blocking reads, bound restart backoff, authenticate local IPC, limit message sizes, and make helper unavailability visible. A helper crash should never preserve motion authority.
\end{advancedbox}

\section{Legacy boundary}

\Legacy\quad Historical controller data uses a fourteen-line keyed payload inside the newer authenticated transport. New domain code should not spread that format. A narrow adapter should translate modern typed commands into the legacy application payload until a deliberate migration replaces it.

A malformed ROBControl header or frame can be rejected before application dispatch. A well-formed, authenticated envelope can still contain a malformed legacy controller payload, so the application adapter must validate that payload separately; transport authentication does not make its contents physically safe.

\section{Typed scene state and local intelligence}
\ROBSourceNow{Cerebro/Cerebro/ROBSceneSnapshot.swift}{Compare the latest-value snapshot store with \texttt{ROBMLXRuntime.swift}, \texttt{ROBMLXImprovisationProvider.swift}, and \texttt{ROBAI.swift}. None of these types is a real-time motor controller.}

\texttt{SceneSnapshot} is the consolidation boundary for people, objects, gestures, depth/lidar free-space regions, arm joint pose, reverse camera pose, camera quality, timestamps, and confidence. Missing observations remain missing. Fresh lidar is retained only briefly; aligned depth is sampled in left/forward/right bands; language serialization labels the JSON as untrusted sensor data rather than instructions. A macOS 26 Foundation Models interpreter may return a typed high-level \texttt{AssistantIntent}, but not servo or tread values.

The local \texttt{ROBMLXEngine} actor owns selected model containers, loading, diagnostics, sampled VLM observations, embeddings, and a process-local semantic store bounded to 200 entries. The inspected initial models are a four-bit Llama 3.2 1B instruct model, a four-bit Qwen2-VL 2B vision-language model, and a small GTE embedding model. The informal term ``VLX'' in project conversation refers to this MLX VLM path, not a separate framework. Camera admission is sampled on a multi-second interval and returns observation text only. It never calls motor code.

Local improvisation must decode as one exact versioned JSON plan with allow-listed beat and delivery values. Gemini remains an optional cloud path with separately effective connection, microphone, camera, and action-tool switches. Both providers produce bounded proposals; deterministic policy, operator authority, and physical executors remain separate.

\ROBPlaceholder{Tag the publication source; provide target macOS and hardware; identify enabled optional dependencies and versions; capture redacted current UI screenshots; document app startup/shutdown, serial-device names, selected Python runtime, and recovery after each subsystem fails.}

\ROBChapterPhoto{Chapter 9}{Controllers and secure transport}{2024-open-tracked-base.jpg}

\chapter{Control from iPhone, Watch, and Vision Pro}
\ROBSourceNow{ROBController/Consciousness/AutoNetClient/AutoNetDataTransferProtocol.swift}{The client control contract is the anchor. Compare it with Cerebro's matching \texttt{AutoNetDataTransferProtocol.swift}; the sending UI is in \texttt{ConsciousViewController.mm}, Watch relay in its Watch app directory, and spatial client under \texttt{ROBControllerVision}.}

\section{Controller roles}

The iPhone/iPad ROBController captures manual controls and presents pairing, autonomy, and action-approval surfaces. Its active legacy snapshot carries raw joystick coordinates; it does not use the Vision controller's normalized linear/angular representation. Apple Watch sends voice or joystick intent through the phone and keeps no robot credential. ROBControllerVision provides a SwiftUI/Swift 6 architecture with a deterministic simulator, normalized motion, a 250 ms lease, a latched software E-stop, secure Cerebro control, and separately transported H.264 video.

\section{ROBControl v2}

\Implemented\quad The production control path advertises \texttt{\_robctl.\_udp}, negotiates \texttt{robctl/2}, uses TLS 1.3 over QUIC, pins Cerebro's certificate, and performs a reciprocal HMAC-SHA-256 proof with a per-device 256-bit secret. Its fixed 40-byte RCTL header carries version, kind, flags, channel, bounded payload length, monotonic sequence, and message ID.

Validation is layered rather than identical for every message. Transport framing checks applicable bounds, role, session, and sequence rules. Lidar telemetry adds capture-time freshness, while controller freshness is measured from Cerebro's local receipt time. The legacy fourteen-line application payload still requires its own decoder validation.

Pairing creates a fresh device UUID and secret for one named device. The server registry remains authoritative for role and revocation. Bonjour advertises non-secret routing facts, not trust. Each device must receive its own credential. The pinned server-certificate fingerprint is a public identity value rather than a secret; security depends on protecting the device secret and private key.

\begin{tabularx}{\textwidth}{>{\robcondensed\bfseries\color{ROBAccent}}p{1.48in} X X}
\toprule
ROLE & MAY & MUST NOT\\
\midrule
operatorController & Send established control/action/autonomy messages and receive results. & Publish lidar as a trusted sensor role.\\
lidarPublisher & Send bounded typed scan and map telemetry. & Request motion, authorize autonomy, or receive generic controller traffic.\\
\bottomrule
\end{tabularx}

\section{Separate video service}

\Implemented\quad Cerebro advertises \texttt{\_robvideo.\_udp} with \texttt{robvideo/1}. The initial profile is H.264 AVCC up to 960 by 540, 20 fps, and 1.5 Mbit/s on an ordered reliable QUIC stream. Video has separate framing, queues, authentication domain, and backpressure so media cannot queue in front of a stop command.

Video authorization is bound to the exact live control-session ID and authenticated controller identity. Control disconnect or credential revocation closes associated video. Video failure does not close control.

\section{Safety timing layers}

\ROBStopChainDiagram

Current timing points include a 250 ms Vision lease, about 450 ms Watch relay watchdog, about 600 ms Cerebro snapshot freshness, and an Arduino loop-count deadline that must be measured. The physical E-stop remains an independent required layer.

\section{Privacy review}

The current active iPhone path requests Always location authorization and includes live latitude and longitude in each controller snapshot. This is current privacy debt, not merely a historical format detail: remove the fields, minimize their precision and retention, or document a necessary purpose and consent policy before public use. Pairing codes, device secrets, private keys, IP addresses, SSH credentials, API keys, private logs, and location must be excluded from screenshots and sample data. A public certificate or fingerprint may be redacted as operational metadata, but it is not a secret. Each cloud-enabled camera or microphone path needs explicit public notice and consent handling.

\ROBPlaceholder{Document exact controller release versions, supported OS/device matrix, pairing and revocation workflow, operator training, latency and packet-loss results, Wi-Fi/router configuration without secrets, public-location data policy, and current Vision Pro hardware validation.}

\ROBChapterPhoto{Chapter 10}{Advanced H.264 video transport}{2022-oakd-depth-camera.jpg}

\chapter{Encode, transport, decode, and present live video}
\ROBSourceNow{Cerebro/Cerebro/ROBCameraH264Encoder.swift}{Follow one frame through \texttt{ROBVideoServer.swift} and \texttt{ROBVideoProtocol.swift}; then cross repositories into ROBControllerVision's \texttt{H264SampleBufferFactory.swift}, \texttt{H264VideoReceiver.swift}, transport decoder, and display-layer bridge.}

This chapter consolidates the robot-side and Vision Pro books into one advanced media path. It describes the implemented initial profile, not a generic H.264 tutorial and not a promise of spatial video, audio, recording, HEVC, or unreliable datagrams.

\section{The end-to-end frame pipeline}

\begin{center}
\begin{tikzpicture}[node distance=0.42cm and 0.24cm,every node/.style={font=\scriptsize}]
\node[draw,rounded corners,fill=ROBCyan!10,minimum width=1.0in,minimum height=0.42in] (capture) {CMSampleBuffer};
\node[draw,rounded corners,fill=ROBAmber!12,right=of capture,minimum width=0.98in,minimum height=0.42in] (scale) {NV12 + letterbox};
\node[draw,rounded corners,fill=ROBAmber!12,right=of scale,minimum width=0.98in,minimum height=0.42in] (encode) {VideoToolbox H.264};
\node[draw,rounded corners,fill=ROBCyan!10,below=of encode,minimum width=0.98in,minimum height=0.42in] (wire) {RVID / RBVD};
\node[draw,rounded corners,fill=ROBCyan!10,left=of wire,minimum width=0.98in,minimum height=0.42in] (validate) {stream validator};
\node[draw,rounded corners,fill=ROBGreen!12,left=of validate,minimum width=1.0in,minimum height=0.42in] (present) {sample renderer};
\draw[->,thick] (capture) -- (scale);
\draw[->,thick] (scale) -- (encode);
\draw[->,thick] (encode) -- (wire);
\draw[->,thick] (wire) -- (validate);
\draw[->,thick] (validate) -- (present);
\end{tikzpicture}
\end{center}

Camera capture produces a Core Media \texttt{CMSampleBuffer} containing a pixel buffer and timing. Cerebro retains only the newest offered raw sample while server work is pending. An accepted, authenticated subscription creates an encoder at the negotiated even width/height, frame rate, and bitrate. The initial published ceiling is 960 by 540, 20 frames/s, and 1.5 Mbit/s.

\section{Scale into the encoder's pixel format}

\texttt{ROBCameraH264Encoder} creates a \texttt{VTCompressionSession} whose destination buffers use video-range bi-planar 4:2:0 YCbCr (NV12). A reusable \texttt{VTPixelTransferSession} scales the source with letterboxing so aspect ratio is preserved. The compression-session pool supplies destination buffers; allocation uses a threshold of six so pressure becomes a dropped frame rather than unbounded buffer growth.

At 20 fps, the target interval is 50 ms. Samples arriving sooner are intentionally skipped before conversion. A dropped allocation or compression result requests a future recovery key frame. Every completion object can be claimed once, preventing synchronous drop flags and an asynchronous VideoToolbox callback from delivering the same logical completion twice.

\section{Configure low-latency H.264}

The encoder requests real-time operation, disables frame reordering, prefers constrained baseline profile when supported, supplies expected frame rate and average bitrate, and limits the key-frame interval to approximately one second. Disabling reordering means this initial stream does not depend on B-frames; decode timestamps can remain invalid and presentation order follows the reliable stream.

Each encoded result owns copied bytes rather than a borrowed VideoToolbox buffer. The encoder records a monotonically increasing access-unit sequence, capture time in Unix milliseconds, presentation time and duration in microseconds, synchronization/key-frame state, NAL-length width, SPS/PPS parameter sets when present, and one complete AVCC access-unit payload.

\section{Understand AVCC, SPS, PPS, and access units}

H.264 consists of Network Abstraction Layer (NAL) units. Annex-B streams separate NAL units with start codes such as \texttt{00 00 00 01}. This implementation does not transmit Annex-B. AVCC stores a big-endian byte length before each NAL unit. The receiver must use the negotiated length-field width; it must not search for start-code patterns inside compressed data.

The Sequence Parameter Set (SPS) describes coded properties needed to construct a decoder format. The Picture Parameter Set (PPS) supplies picture-level decoding parameters. Cerebro sends them in a codec-configuration message with a positive generation before access units using that generation. A new configuration, encoder recovery, or requested key frame resends the configuration. An IDR key frame lets the decoder restart without depending on earlier reference pictures.

\begin{advancedbox}
Codec configuration is state, not decoration. A payload can be valid H.264 bytes yet undecodable with the wrong SPS/PPS, NAL-length width, dimensions, stream ID, or generation. Validate all of them before asking Core Media to allocate decoder state.
\end{advancedbox}

\section{Wrap media in two bounded layers}

The outer ordered-stream frame is \texttt{RVID} version 1. Its 32-byte header carries magic, version, header length, message kind, bounded payload length, monotonic connection sequence, and reserved-zero fields. Authentication, capability, subscription, feedback, configuration, access-unit, and end messages receive distinct kinds.

Encoded configuration and access units use an inner \texttt{RBVD} version-1 header of 92 bytes. It binds media to session UUID, subscription UUID, media sequence, capture time, presentation time, duration, timescale, codec-configuration generation, NAL-length width, flags, and payload length. Integers are network byte order. Current hard payload limits are 64 KiB for codec configuration and 2 MiB for one access unit; control JSON is capped at 64 KiB and rejects unknown fields.

\begin{tabularx}{\textwidth}{>{\robcondensed\bfseries\color{ROBAccent}}p{1.2in} p{1.25in} X}
\toprule
BOUNDARY & PRIMARY CHECKS & WHY IT EXISTS\\
\midrule
TLS/QUIC & pinned identity, proof, role & Authentic peer and protected ordered bytes.\\
RVID & magic, version, kind, length, sequence & Incremental connection framing before allocation.\\
RBVD & IDs, codec, timing, generation, lengths & Media belongs to the negotiated live stream.\\
AVCC & NAL lengths and configuration & Compressed payload can be reconstructed safely.\\
Core Media & coded/presentation dimensions & SPS/PPS cannot silently exceed negotiated geometry.\\
\bottomrule
\end{tabularx}

\section{Apply backpressure before latency accumulates}

Video has its own listener, connections, queue, encoder, authentication domain, and send state. It never shares the control message queue. Cerebro permits only one pending raw source frame, one pending encoder output per subscription, and one media send in flight. If encoding or send work is occupied, newer unencoded work is dropped rather than queued. One actual encoded drop or undeliverable output forces a later IDR; replacing a not-yet-encoded raw frame does not break decoder references.

Every RVID send has a ten-second completion deadline. A newly accepted stream has 15 seconds to produce its first frame, and a running stream may not remain silent beyond 15 seconds. A peer that stops reading loses the stream instead of retaining unlimited memory. At the published maximum, two viewers can still consume roughly 3 Mbit/s plus protocol overhead, so data-plane separation prevents queue coupling but not shared-radio contention.

\section{Decode framing before decoding H.264}

On Vision Pro, the Network.framework adapter first reconstructs complete RVID frames from arbitrary receive chunks. It rejects invalid magic, version, reserved bytes, payload sizes, and ordering before decoding JSON or RBVD. The media decoder then checks session ID, subscription ID, codec, configuration generation, sequence, timing, AVCC length width, and size. These checks occur before creating Core Media state.

\texttt{H264SampleBufferFactory} sorts and extracts SPS/PPS, copies their bytes into temporary stable pointers, and calls \texttt{CMVideoFormatDescriptionCreateFromH264ParameterSets}. It verifies that coded dimensions remain within macroblock rounding of the negotiated size and that clean-aperture presentation dimensions match. This prevents hostile or corrupted parameter sets from expanding the decoder beyond the agreed geometry.

For each accepted access unit, the factory allocates a \texttt{CMBlockBuffer}, copies the bounded AVCC payload, creates timing, and asks Core Media to create a ready sample buffer. The not-sync attachment mirrors the key-frame flag. The display-immediately attachment requests low-latency presentation rather than building a long presentation queue.

\section{Present and recover without blocking control}

The actor-owned \texttt{H264VideoReceiver} feeds an \texttt{AVSampleBufferVideoRenderer} obtained from the UI's display layer. If the renderer is temporarily full, it waits only about 50 ms in short cancellable intervals. Continued pressure drops the access unit and begins recovery rather than accumulating frames.

Sequence gaps, missing configuration, generation mismatch, decoder failure, or \texttt{requiresFlushToResumeDecoding} make the validator require a key frame. Recovery flushes decoder presentation state, rate-limits key-frame feedback to no faster than 500 ms, and discards dependent frames until a valid key frame arrives. A configuration-generation change also flushes before accepting the new decoder state. Statistics report received, rendered, dropped, bytes, last sequence, and key-frame requests at a bounded update rate.

\section*{Timing, reliability, and tradeoffs}

Capture wall-clock time measures end-to-end age. Presentation timestamps provide an internal monotonic media timeline. Duration is derived from the negotiated frame rate. These clocks answer different questions and should not be substituted for one another.

Ordered reliable QUIC simplifies configuration and access-unit ordering, but packet loss can delay later video bytes through head-of-line blocking. Because control uses another QUIC connection, video loss does not queue a stop command, although both flows share Wi-Fi airtime. Future QUIC datagrams would trade retransmission for application-managed loss recovery and require a new negotiated contract; they are not silently enabled by the current enum.

\section{Advanced verification matrix}

Test the media pipeline with synthetic frames before live hardware:

\begin{itemize}
  \item verify scaling, even dimensions, NV12 allocation limits, bitrate, frame-rate skipping, IDR interval, and SPS/PPS extraction;
  \item reject truncated, oversized, mis-versioned, wrong-session, wrong-generation, unknown-field, and out-of-order messages before allocation;
  \item inject access-unit loss and confirm one bounded recovery sequence and key-frame request;
  \item change codec generation and confirm flush, reconfiguration, and key-frame resume;
  \item stall the renderer and network independently and prove memory remains bounded;
  \item close or revoke the control session and prove its video subscription ends;
  \item measure capture-to-display age, encoded size, drops, recovery time, CPU/GPU use, thermal state, and control latency concurrently;
  \item repeat camera unplug, app background, viewer close, and shutdown while a send and encode callback are in flight.
\end{itemize}

Volume 7 explains the Swift/visionOS client composition and user-facing display. Volume 8 explains Cerebro's broader camera and model scheduling. This chapter is the shared wire-level and Core Media reference.

\ROBChapterPhoto{Chapter 11}{Cameras, depth, lidar, and coordinate frames}{2022-oakd-depth-camera.jpg}

\chapter{Perception and mapping}
\ROBSourceNow{Cerebro/Cerebro/CameraManager.swift}{Camera ownership starts here. Depth helpers remain in Cerebro; lidar transport and fixtures live in the separate RPLidar repository named in the repository source map.}

\section{RGB-D camera}

\Experimental\quad The current camera design runs a pinned DepthAI helper and sends timestamp-synchronized RGB plus depth through a private local Unix socket. RGB enters the existing sample-buffer, Vision, Gemini, and video path. Depth is little-endian UInt16 millimeters aligned to RGB; zero is invalid. AVFoundation remains an RGB-only fallback.

The helper owns the OAK device exclusively. Old frames are dropped when perception is busy. Missing hardware, malformed IPC, SDK failure, or helper exit becomes camera state and bounded retry rather than a control-process crash.

The helper uses DepthAI Camera, Depth, and Sync nodes, requests 640 by 400 undistorted RGB, aligns RAW16 millimeter depth to RGB, and emits a versioned \texttt{CDP1} packet over a private Unix socket. Cerebro validates dimensions, pixel formats, byte counts, and maximum allocation before constructing a CoreVideo buffer. The process boundary contains SDK and USB failure; a future native or XPC helper should preserve the same \texttt{CameraFrameSet} provider contract.

\section{Perception snapshots and bounded rendering}

Camera rate is not inference rate. Preview may consume the newest RGB frame while Vision pose detection, colored depth, point-cloud geometry, Gemini JPEG sampling, MLX VLM sampling, and H.264 encoding run at distinct bounded rates. Slow consumers use latest-value mailboxes. SceneKit reuses or replaces named node roots instead of appending forever; Core Image contexts and pixel-buffer pools are reused; model containers stay actor-owned.

Profile repeated source toggles and 30-minute runs, not one still frame. Look for retained \texttt{CMSampleBuffer} objects, per-frame \texttt{CIContext} construction, ever-growing SceneKit textures or nodes, missing autorelease pools on Objective-C worker threads, duplicate model loads, and repeated diagnostics overwhelming the main thread.

\section{Kinect, OpenNI, NiTE, and PCL archaeology}

\Legacy\quad The \texttt{ROBNiTEManager}, \texttt{TaskControllers}, \texttt{FreenectPCL}, \texttt{k2g}, viewer, shader, serialization, and duplicated ``2'' files preserve an older Kinect/OpenNI/NiTE/PCL depth stack. They entered the present Git history in the 2025 initial import, which does not establish their original creation date. They are useful evidence for coordinate and visualization experiments, not the recommended OAK setup.

Before removal, check Xcode target membership, runtime selectors, duplicated model provenance, and any coordinate knowledge that should move into tests or documentation. The modern DepthAI helper improves failure containment by placing the camera SDK and USB ownership outside the control process.

\section{RPLidar}

\Experimental\quad The companion app reads pose, location, laser points, and occupancy maps from Slamware. It publishes scans about every 0.2 seconds and maps about every 5 seconds through typed, sequenced, bounded telemetry. Before scan publication, the current code keeps only valid points closer than \(1\,\mathrm{m}\) and within approximately \(\pm0.50\) radians of the front or back; this scan is not a full 360-degree view. The occupancy map is a separate payload. The publisher role cannot drive the robot.

Cerebro checks authenticated publisher identity, monotonic sequence, capture time, kind, and bounds. Stale data is rejected. Current social-roam autonomy consumes fresh scans, not occupancy maps.

\section{IR and IMU}

\Experimental\quad Front/back Sharp IR warnings remain active, but automatic movement-inhibit flag assignments are disabled. The floor/grass reflection history makes these sensors an ideal research topic, not a certified stop device.\par

\Implemented\quad The MPU9250 reports inertial data and derived orientation. It is not a motion-inhibit input.

\section{Coordinate-frame discipline}

Name every transform:

\begin{itemize}
  \item lidar sensor to base;
  \item base to local map;
  \item color camera to depth camera alignment;
  \item camera to head/neck;
  \item head/neck to torso;
  \item torso to each arm base;
  \item arm joint chain to gripper tool frame.
\end{itemize}

Record axis direction, handedness, units, timestamp, calibration method, uncertainty, and valid motion range. A visually plausible transform can still command a gripper to the wrong side of an object.

\ROBPhotoTwo{rob-v3-page-55.jpg}{Presentation page 55 proposes an RPLidar stage above the base.}{2024-open-tracked-base.jpg}{The open base shows how sensors, wiring, compute, and mechanisms become a serviceability problem.}

\ROBPlaceholder{Confirm current camera models, lenses, mounting transforms, USB/Ethernet topology, DepthAI version, exposure and lighting behavior, lidar model/host/address, installed map workflow, IR positions and shrouds, calibration targets, timestamp synchronization, frame conventions, and measured accuracy.}

\ROBChapterPhoto{Chapter 12}{Arms, animation, speech, autonomy, and AI}{2025-detached-articulated-arm.jpg}

\chapter{Behavior above raw motor commands}
\ROBSourceNow{Cerebro/Amber-PythonAPI/Amber V2 API/amber_api/amber_robot.py}{This is the readable high-level AMBER client. Compare its packed command files with the Ubuntu launch JSON and LCM schemas; the captured \texttt{amber\_core} executable is not open source in this snapshot.}

\section{Servo and arm control}

Cerebro uses Maestro servo channels plus AMBER Python paths. ROB's present AMBER evidence describes two seven-joint B1 arms and separate grippers. The Ubuntu-side directories \texttt{L-10} and \texttt{R-11} bind \texttt{can10}/\texttt{can11}, UDP ports 26001/26002, and LCM prefixes \texttt{Left\_}/\texttt{Right\_}. Network clients send UDP or LCM to one core per side; each core, not the client, interfaces with its actuators over CAN.

The checked-in V2 protocol provides status, joint and Cartesian moves, calibration, gripper control, and mode control. Its packed UDP header contains a 16-bit command, 16-bit length, and 32-bit counter. The current client needs stronger validation of source, size, command, counter, freshness, finite values, and authentication before production use. Archived binaries do not include enough source to reconstruct their actuator-level CAN protocol.

URDF models provide link geometry, joint origins, axes, and limits, but they do not prove ROB's physical mount transforms, encoder zero, payload, safe stopping, or cable collision volume. Measure torso-to-left-base, torso-to-right-base, and each tool transform; then verify forward kinematics over several supported poses before inverse kinematics. Volume 6 contains the protocol tables, joint model, Ubuntu reproduction procedure, and staged dual-arm commissioning gates.

\section{Keyframe animation}

\Implemented\quad \texttt{KeyframeAnimationManager.swift} persists JSON values for arm channels, head/neck, tread movement, flipper, actuator, torso rotation, and speech. Keyframes are reusable data: a sequence can describe expressive behavior independently of the UI that recorded it.

\Commanded\quad Historical servo ``animation'' helpers in \texttt{ROBSerialBox.m} do not yet perform robust smooth interpolation; they update their targets immediately. New interpolation must respect joint rates, synchronization, limits, collisions, and cancellation.

\section{Speech}

\Implemented\quad Cerebro configures an Apple local speech-recognition path and uses local speech synthesis for completed responses. This path reduces network dependence, but it does not describe every microphone use: when the separately optional Gemini session is enabled, microphone audio may be streamed to an external cloud service. Both paths require microphone permissions, public notice, and a plan for unintended capture.

\section{Controller-activated autonomy}

\Experimental\quad The social-roam profile begins only after explicit controller authorization, captures an activation pose and radius, uses fresh filtered lidar scans, and limits tread speed. Near 82 percent of the configured radius it attempts to return toward its origin; this is not verified containment. Missing or stale lidar requests a motion stop and transitions to \texttt{waitingForLidar} while the autonomy session remains active. Manual takeover, explicit stop, expiry, replacement, or shutdown ends the session. Local deterministic code owns execution.

\section{Gemini integration}

\Experimental\quad The entire Gemini integration defaults off and requires explicit enablement plus credentials. Once enabled, the current configuration defaults camera-frame and microphone-audio streaming on unless each stream is separately disabled. The robot-action tool remains separately off by default. Its schema declares \texttt{look\_at}, \texttt{play\_gesture}, \texttt{request\_pick}, \texttt{navigate\_relative}, and \texttt{stop\_motion}; the model never receives arbitrary raw serial output. Only \texttt{stop\_motion} has an immediate local executor, and even that returns a partial result because arm hold is unverified. The other actions have no deterministic executor: the controller console records approval or rejection but does not actuate hardware.

\Planned\quad General picking is intentionally unavailable. Missing prerequisites include calibrated camera-to-arm transforms, complete neck mapping, joint feedback, inverse kinematics, collision checking, grasp planning, and a verified executor.

\section{Local MLX, comedy, and battle training}

\ROBSourceNow{Cerebro/Cerebro/ROBStageShowCoordinator.swift}{Read the strict show protocol and fixture tests first, then compare \texttt{ROBMLXImprovisationProvider.swift}, \texttt{ROBSaberChoreography.swift}, and the JSON catalog under \texttt{StageShows}.}

The stage-show coordinator is a cancellable state machine rather than a sequence of arbitrary shell or motor commands. Show JSON may request speech, waits, checkpoints, allow-listed gestures, and adaptive dialogue; it may not contain raw joints, hosts, ports, or shell commands. Dry Run has no side effects. Offline uses authored lines. Local mode may use a schema-validated MLX or loopback llama.cpp line. Adaptive mode may give a trusted local plan to Gemini, but timeout and failure fall back through validated local text to mandatory authored text. Late model completions are correlated and discarded.

The comedy show demonstrates why fallback belongs to creative design: roughly ten minutes of timing cannot depend on one network or model response. Progressive saber training maps named profiles to bounded Cartesian transforms and duration. It must remain theatrical, slow, guarded choreography with soft props, a participant boundary, dry-run visualization, independent stop, transform, inverse-kinematics, and collision validation, and no autonomous contact. Vision may score timing or offer spoken coaching; pixels do not directly command the AMBER arms.

\section{Robot-action implementation status}

\begin{tabularx}{\textwidth}{>{\ttfamily\scriptsize}p{1.35in} p{1.35in} X}
\toprule
ACTION & CURRENT STATUS & CURRENT BEHAVIOR\\
\midrule
stop\_motion & Partial local executor & Preempts local coordinators, writes one neutral/braked base frame, and drops the base heartbeat. Base stop is software-observed; Amber arm hold remains unverified.\\
look\_at & Schema only & No deterministic local executor. Outside autonomy it may be offered to the manual approval console; approval does not move hardware.\\
play\_gesture & Schema only & No deterministic local executor. Approval does not play a physical gesture.\\
navigate\_relative & Schema only & No deterministic local executor. It is distinct from the social-roam coordinator.\\
request\_pick & Planned / unavailable & Rejected until transforms, inverse kinematics, collision checking, joint feedback, grasp planning, and a verified executor exist.\\
\bottomrule
\end{tabularx}

\begin{safetybox}
Language-model confidence is not evidence of physical safety. Treat model output as an untrusted proposal. Validate types, values, state, authority, freshness, workspace, and collision policy in deterministic local code; require human approval where risk warrants it; and make rejection an ordinary result.
\end{safetybox}

\ROBPhotoTwo{2025-detached-articulated-arm.jpg}{The detached arm makes reach, joint order, cabling, and manipulation limits concrete.}{2022-oakd-depth-camera.jpg}{The depth camera is one input to a manipulation experiment; a successful image alone is not a grasp-validation test.}

\ROBPlaceholder{Provide arm/neck calibration files, verified keyframes, joint feedback and limits, safe animation cancel behavior, speech privacy settings, autonomy speed/radius/obstacle requirements, Gemini model/release configuration, consent signage, action approval policy, and validation for every enabled tool.}

\part{Integration, operation, and maintenance}

\ROBChapterPhoto{Chapter 13}{Verification before motion}{2022-tracked-base-front.jpg}

\chapter{Test strategy and safety case}

\section{Requirements traceability}

Every test should trace to a requirement and hazard. A simple matrix includes requirement ID, hazard ID, source revision, test method, equipment, pass criteria, result, evidence path, reviewer, and date. ``It worked once'' is not a repeatable acceptance criterion.

\section{Progressive test ladder}

\begin{tabularx}{\textwidth}{>{\robcondensed\bfseries\color{ROBAccent}}p{0.6in} X X}
\toprule
LEVEL & SCOPE & REQUIRED EVIDENCE\\
\midrule
0 & Design review and hazard analysis. & Reviewed drawings, calculations, assumptions, and unresolved risks.\\
1 & Pure software fixtures and simulation. & Protocol, range, state, deadline, replay, and cleanup tests.\\
2 & Logic outputs to indicators or dummy loads. & Pin states, ordering, safe initialization, timeout waveform.\\
3 & One controller with a small surrogate or disconnected actuator. & Polarity, limits, current, fault and stop response.\\
4 & Full hardware on rated stands with guards. & Direction, speed cap, coast/brake, current, temperature, E-stop timing.\\
5 & Slow unloaded floor test in an exclusion zone. & Tracking, latency, stopping distance, spotter log, post-test inspection.\\
6 & Incremental loads and declared surfaces. & Repeatable operating envelope and thermal/current margins.\\
7 & Full show rehearsal. & Setup checklist, barrier, staffing, failure drills, privacy notice, teardown plan.\\
\bottomrule
\end{tabularx}

\section{Failure injection}

Test the robot's response when:

\begin{itemize}
  \item joystick input stops or becomes non-finite;
  \item Watch-to-phone relay disappears;
  \item controller packets duplicate, reorder, age, or exceed bounds;
  \item control network disappears while video continues;
  \item video decoder stalls or requests a keyframe;
  \item USB serial is unplugged or filled with malformed bytes;
  \item Arduino boots with IMU missing;
  \item an IR sensor sticks low or produces grass reflections;
  \item lidar data becomes stale or claims an impossible timestamp;
  \item DepthAI helper exits or the camera disconnects;
  \item an actuator reaches a mechanical stop;
  \item one motor direction is reversed;
  \item a credential is revoked during active control;
  \item autonomy loses its sensor or leaves its radius;
  \item an AI tool proposes an unsupported or stale action;
  \item physical E-stop is pressed at each declared operating speed and pose.
\end{itemize}

\section{Measure stop behavior}

For each layer, measure detection time, command time, electrical response, mechanical response, and final state. Report surface, slope, load, battery state, speed, controller mode, firmware/software revision, sample rate, and uncertainty. ``600 ms watchdog'' is not the same as ``stopped within 600 ms.''

\begin{safetybox}
Use rated stands, restraints, remote observation, current limiting, a spotter, and an independent stop before intentionally injecting failures. Some failures should be tested only in simulation or on surrogate hardware.
\end{safetybox}

\ROBPlaceholder{Insert the actual hazard analysis, operating envelope, requirements matrix, unit and fixture reports, current/thermal data, watchdog timing, E-stop test, stopping distances on hard floor and grass, slope tests, flipper/actuator load tests, arm collision tests, and sign-off authority.}

\ROBChapterPhoto{Chapter 14}{Maker Faire operations}{2025-rob-side-full-body.jpg}

\chapter{Present ROB safely in public}

\section{Roles}

Assign at least these responsibilities:

\begin{itemize}
  \item \textbf{Lead operator}: owns arm/disarm and the primary controller.
  \item \textbf{Safety spotter}: watches the robot and crowd, keeps a hand near the physical E-stop, and may stop without asking.
  \item \textbf{Educator}: explains ROB without distracting the operator.
  \item \textbf{Technical support}: monitors power, temperature, logs, and recovery.
  \item \textbf{Event liaison}: coordinates space, charging, barriers, accessibility, and emergency requirements.
\end{itemize}

One person may fill multiple roles only if workload and risk allow it. Never make the operator answer a long technical question while ROB moves near visitors.

\section{Pre-show checklist}

\begin{itemize}[label={\protect\raisebox{0.08em}{\protect\fbox{\protect\rule{0pt}{0.55em}\protect\rule{0.55em}{0pt}}}}]
  \item current approved software/firmware versions recorded;
  \item batteries inspected, charged under the approved procedure, and restrained;
  \item fuses, disconnect, and E-stop visually inspected and function-tested;
  \item treads, chain, sprockets, guards, shafts, collars, mounts, fasteners, and torque marks inspected;
  \item arm, neck, flipper, and actuator limits checked in a protected area;
  \item controllers uniquely paired; unused credentials revoked;
  \item control and optional video links tested under venue network conditions;
  \item speed cap and operating mode verified on every controller;
  \item barrier, route, slope, floor transitions, cables, exits, and charging location inspected;
  \item camera, microphone, recording, location, and cloud-AI policy displayed;
  \item fire response, first aid, lift/transport, and fault-recovery plans briefed;
  \item spare tools and parts controlled so visitors cannot reach them;
  \item dry rehearsal completed with the full team.
\end{itemize}

\section{Show-state model}

Use explicit visible states: Transport Safe, Display Only, Demo Armed, Moving, Stopped, Faulted, Charging, and Maintenance. Each transition names who authorizes it and what checklist applies. Reconnecting a controller or restarting software never silently arms motion.

\section{Audience and privacy}

Keep an accessible viewing route. Do not require visitors to step over cables. Mark pinch and tread boundaries clearly. Avoid publishing faces without consent. When camera, microphone, recording, or cloud processing is active, use plain-language signage and a visible indicator; provide a quick disable path and a non-recorded way to participate.

\ROBPlaceholder{Insert venue-specific barrier dimensions, maximum public-demo speed, operator and spotter names, E-stop placement, battery/charging policy, transport vehicle and lift plan, insurance or venue requirements, photo/video release plan, privacy signage, scripted demonstrations, and abort criteria.}

\ROBChapterPhoto{Chapter 15}{Maintenance and change control}{2025-rear-electronics-bay.jpg}

\chapter{Keep the robot understandable}

\section{Before and after every operating day}

Inspect for loose fasteners, moved torque marks, chain and tread damage, shaft-collar movement, bracket cracks or bending, wire abrasion, connector heating, battery swelling or damage, fluid or debris, blocked ventilation, actuator looseness, unusual sound, camera security, E-stop function, and log errors.

After operation, record runtime, surfaces, mode, maximum observed current and temperature if available, incidents, unexpected stops, repairs, and deferred defects. Quarantine the robot from powered use when a red condition remains unresolved.

\section{Scheduled maintenance record}

For every interval, name the component, task, acceptable condition, tool, replacement part, torque or measurement, evidence, and sign-off. Intervals may be hours, cycles, calendar time, or before each public show.

\section{Configuration control}

A change to code, pin mapping, motor polarity, brake logic, network framing, controller speed cap, arm limit, battery arrangement, or bracket can invalidate earlier tests. Give hardware, firmware, software, calibration, wiring, and operating procedure separate revision identifiers, then record the tested combination.

\section{Troubleshooting by layer}

\begin{tabularx}{\textwidth}{>{\robcondensed\bfseries\color{ROBAccent}}p{1.35in} X}
\toprule
SYMPTOM & SAFE DIAGNOSTIC PATH\\
\midrule
No controller connection & Keep motion disarmed; verify unique pairing, robot ID, Bonjour visibility, certificate pin, role, revocation, and session proof.\\
Control works, no video & Keep control independent; inspect the separate video service, camera permission/state, capability handshake, session binding, and decoder recovery.\\
Arduino not online & Disconnect actuator power as required; verify USB device, 250,000 baud, startup log, safe GPIO state, and dependency/IMU behavior.\\
One tread reverses & Stop and lift/guard the base; compare requested sign, adapter output, Arduino pin, driver polarity, and physical mount. Do not fix by guessing in UI code.\\
Unexpected IR warnings & Disable motion; log all raw readings with surface, angle, light, and distance; inspect shrouds and wiring; keep IR outside certified safety claims.\\
Autonomy will not move & Treat as safe failure; verify explicit authorization, fresh lidar, activation radius, master control, and local planner state.\\
Arm action rejected & Treat rejection as correct until capability, calibration, approval, and executor prerequisites are satisfied.\\
\bottomrule
\end{tabularx}

\ROBPlaceholder{Add maintenance intervals, lubrication policy, tread and chain wear limits, fastener torques, battery retirement criteria, cable replacement rules, calibration intervals, spare-parts inventory, fault codes, backup/restore procedure, and configuration-release form.}

\appendix

\chapter{Teach the architecture without hiding the hazards}

The youth sequence and public website should use the same conceptual spine as the engineering record: system boundary; energy versus information; sense--estimate--plan--validate--act--measure; mechanics and operating envelope; message grammar and freshness; identity and authority; progressive evidence. The explanation may become simpler, but it must not become false.

\section{A spiral curriculum for ROB}

Return to the same ideas with increasing precision. Ages 8--12 classify inputs, outputs, energy, structure, and feedback using paper models. Ages 10--14 graph signals, parse bounded messages, and implement current-limited LED watchdogs. Ages 10--15 model forces, ratios, tolerances, and test ladders. Ages 12--16 reason about state machines, networking, trust, leases, perception limits, and bounded autonomy. Adult mentors connect each model to dated evidence and unresolved hazards.

\begin{tabularx}{\textwidth}{>{\robcondensed\bfseries\color{ROBAccent}}p{1.15in} X X}
\toprule
LEARNING MOVE & YOUTH EVIDENCE & ENGINEERING EVIDENCE\\
\midrule
Observe & annotated photo or sensor-map marks & calibrated log with timestamp and revision\\
Model & block diagram, graph, state cards & interface control document or executable model\\
Predict & written normal/change/failure result & acceptance criterion with uncertainty\\
Test & paper, simulation, or low-energy kit & progressive restrained test with hazard controls\\
Explain & claim linked to observation & traceable result, exception, and sign-off\\
Revise & changed diagram and reason & controlled design revision and regression test\\
\bottomrule
\end{tabularx}

\section{Facilitation pattern}

Begin with a phenomenon, not a vocabulary list. Ask learners to notice; collect competing models without ridicule; choose a measurable prediction; test at the lowest safe energy; compare result and prediction; then introduce precise language. End by asking what remains unknown. This preserves curiosity while teaching that uncertainty is an engineering result.

\begin{safetybox}
Do not turn ROB into a youth build activity. The public curriculum must not invite access to traction batteries, mains or inverter circuits, drivetrain hardware, chains, machine tools, actuators, live controllers, or unrestricted network commands. Demonstrations use immobilized hardware, simulation, paper, or separately reviewed low-energy kits.
\end{safetybox}

\chapter{Commission the control path as one safety-related system}

The transport should be verified from discovery through physical response. The current secure path uses Bonjour discovery, QUIC/TLS, exact certificate pinning, reciprocal pairing proof, a server-owned device role, framed application messages, controller freshness, Cerebro's base heartbeat ownership, and the Arduino's independent serial timeout. Each layer needs its own test because success at one layer cannot establish another.

\section{Connection and enrollment matrix}

Record clean enrollment, reconnect after launch-order reversal, local-network permission denial, missing Bonjour TXT metadata, wrong certificate, wrong secret, revoked device, duplicated credential, wrong role, stale sequence, Wi-Fi loss, and server restart. Capture both endpoint states and verify that no application callback occurs before authentication completes.

\begin{fieldnotebox}[title={2026-08-10 TRANSPORT LESSON}]
During live commissioning, iOS discovered the sole \texttt{\_robctl.\_udp} service but did not surface its Bonjour TXT record. The controller now permits a connection attempt only for that sole metadata-less candidate; the exact certificate pin and reciprocal proof still fail closed for another robot. A second issue exposed QUIC anti-amplification behavior: Cerebro attempted to send the challenge before the client address had enough validation traffic. The protocol now begins with a padded client pairing-hello carrying a registry-checked device identifier, then sends the fresh challenge. Treat this as a dated implementation note pending tagged release and regression fixtures.
\end{fieldnotebox}

\section{Measure the whole stop chain}

Instrument controller issue time, validated receipt, Cerebro neutralization, last serial frame, driver output, tread speed, and complete physical stop. Inject joystick release, app suspension, cable removal, malformed frames, serial noise, stale sensors, helper failure, and E-stop activation separately. Report both software state and measured motion. A green authenticated-link indicator proves identity and readiness; it does not prove command freshness, brake behavior, or stopping distance.

\begin{advancedbox}[title={RELEASE ARTIFACTS}]
Archive source commits, pairing protocol version, component configuration, redacted logs, fixture results, hazard review, calibration set, operator checklist, known limitations, and rollback procedure. Make the public book and web lesson point to a release record rather than to a developer's working tree.
\end{advancedbox}

\chapter{Capture the builder's oral history}

The future first-person story should be dictated as evidence-rich scenes, not forced into a polished chronology in one sitting. Record separate sessions for origin, early moving base, body and character, power and electronics, fabrication, sensing and software, public demonstrations, failures, collaborators, and the present mission. Preserve the speaker's exact recording privately; produce a reviewed transcript; mark uncertain dates; and link claims to photographs, commits, drawings, or test records.

\section{Interview prompts}

For each revision ask: What problem were you trying to solve? What did you believe then? What did you build? What surprised or failed? What changed next? Who helped? Which artifact proves the memory? What would you tell a child, and what would you tell an experienced builder? Do not fill technical placeholders from recollection when a measurement is required; use the story to identify the measurement or archive item to find.

\begin{missionbox}[title={DICTATION WORKFLOW FOR THE NEXT EDITION}]
Record one 20--30 minute chapter at a time. Start each with date range and visible artifact. Transcribe, separate direct memory from later interpretation, flag names and permissions, fact-check technical claims, then return the edited passage to the builder for voice and accuracy approval. Only approved passages enter the narrative edition.
\end{missionbox}



\chapter{Builder input ledger}

This appendix is the iteration queue. Replace each line with a dated, reviewed answer and evidence reference.

\begin{enumerate}
  \item Project origin, name meaning, design requirements, contributors, shops, and dated revision history.
  \item Final overall dimensions, ready-to-run mass, center of mass, load cases, operating surfaces, slopes, speed, and environmental limits.
  \item As-built CAD/templates, plate materials and thicknesses, fasteners, shafts, bearings, sprockets, chains, tension procedure, guards, and coatings.
  \item Actual milling, drilling, bending, welding, deburring, finishing, workholding, tooling, feeds/speeds, inspection, PPE, and vendor sequence.
  \item Current tread and flipper motors, windings, gearbox ratios, controller circuits, brake polarity, current limits, and performance tests.
  \item Flipper location and geometry, third-wheel status, actuator roles, models, strokes, loads, end stops, feedback, and failure behavior.
  \item Battery packs, topology, BMS, charger, fusing, main disconnect, E-stop, converters, inverter, grounding, wire gauges, connectors, and power budget.
  \item Arduino-compatible board model, core and libraries, pin-to-connector harness, command ranges, deadman timing, IMU calibration, IR positions/status, and speed-feedback status.
  \item Current Mac mini/Raspberry Pi inventory, OS versions, ports, network devices, router role, serial-device ownership, and dependency versions.
  \item Arm and neck models, geometry, calibration, zero/home, soft/hard limits, payload, current/torque limits, collision envelope, and safe transport.
  \item Camera and lidar inventory, mounts, coordinate transforms, calibration, timestamp synchronization, accuracy, lighting, and present host path.
  \item Controller compatibility, pairing/revocation, speed caps, latency, watchdogs, location-data policy, and Vision Pro validation.
  \item Autonomy operating envelope, sensor requirements, aborts, AI model/tool configuration, approval policy, and camera/microphone/cloud consent.
  \item Hazard analysis, protective measures, E-stop proof, test reports, Maker Faire perimeter and staffing, fire, lifting, and charging plans, plus insurance and venue requirements.
  \item Final photo captions, permissions/releases, vendor-publication permissions, credits, and redacted screenshots.
\end{enumerate}

\chapter{Bill of materials worksheet}

\begin{tabularx}{\textwidth}{p{0.35in} p{1.15in} X p{0.55in} p{0.75in} p{0.8in}}
\toprule
ID & SUBSYSTEM & MANUFACTURER / PART / REVISION & QTY & SOURCE & STATUS\\
\midrule
\rule{0pt}{0.34in} & & & & &\\
\rule{0pt}{0.34in} & & & & &\\
\rule{0pt}{0.34in} & & & & &\\
\rule{0pt}{0.34in} & & & & &\\
\rule{0pt}{0.34in} & & & & &\\
\rule{0pt}{0.34in} & & & & &\\
\rule{0pt}{0.34in} & & & & &\\
\rule{0pt}{0.34in} & & & & &\\
\bottomrule
\end{tabularx}

For each line add voltage, current, mechanical interface, firmware, calibration, supplier status, substitute, unit cost, lead time, safety documentation, and evidence photograph in the editable project record.

\chapter{Wiring and connector worksheet}

\begin{tabularx}{\textwidth}{p{0.55in} p{1in} p{1in} p{0.65in} p{0.58in} p{0.7in} X}
\toprule
NET & FROM & TO & VOLTS & MAX A & WIRE & PROTECTION / CONNECTOR\\
\midrule
\rule{0pt}{0.34in} & & & & & &\\
\rule{0pt}{0.34in} & & & & & &\\
\rule{0pt}{0.34in} & & & & & &\\
\rule{0pt}{0.34in} & & & & & &\\
\rule{0pt}{0.34in} & & & & & &\\
\rule{0pt}{0.34in} & & & & & &\\
\bottomrule
\end{tabularx}

Also record return path, chassis bond, shield termination, cable length, routing class, strain relief, flex cycles, color/label, fuse/breaker location, test point, polarity check, and continuity/insulation result.

\chapter{Serial protocol reference}

\section*{Historical base frame}

\begin{tabularx}{\textwidth}{>{\robcondensed\bfseries}p{0.65in} p{1.55in} p{1.1in} X}
\toprule
FIELD & MEANING & EXAMPLE & CURRENT LIMIT / NOTE\\
\midrule
Start & Frame marker & \texttt{\textasciitilde} & One byte\\
1 & Left brake & \texttt{+0000} & 0 described as release; nonzero engage\\
2 & Left signed speed & \texttt{+0100} & Expected magnitude within 0-255; not clamped\\
3 & Right brake & \texttt{+0000} & Same mapping\\
4 & Right signed speed & \texttt{+0100} & Mirrored physical direction polarity\\
5 & Flipper brake & \texttt{+0000} & Same mapping\\
6 & Flipper signed speed & \texttt{+0000} & Expected magnitude within 0-255; not clamped\\
7 & Actuator signed speed & \texttt{+0000} & Nominal -3200 to +3200; +3201 and -3201 clear safe start, then command +3200 and -3200 respectively\\
\bottomrule
\end{tabularx}

Fields are five ASCII characters and separated by six commas. Total historical frame length is 42 bytes. The current decoder defects listed in Chapter 7 must be fixed before treating the format as safe. A replacement should add explicit version, message type, bounded length, sequence, integrity check, and typed state transitions.

\chapter{Repository source map}

\begin{tabularx}{\textwidth}{>{\ttfamily\scriptsize}p{2.65in} X}
\toprule
RELATIVE PATH & PUBLICATION USE\\
\midrule
Presentation/ROB\_v3.pdf & 95-page historical build and design deck; slides dated 2021 with 2022 modifications.\\
ROBArduino/ROBOT\_CEREBELLULAR\_BASE\_APP/...ino & Present Base-role reference: pin map, parser, treads, flipper, actuator, IR, IMU, heartbeat.\\
ROBArduino/ROBOT\_CEREBELLULAR\_TORSO\_APP/...ino & Retired Torso-role history: Maestro head/arm targets, IMU, and disabled GPS/timeout behavior.\\
ROBArduino/ROBOT\_CEREBELLULAR\_HEAD\_APP/...ino & Retired Head-role history: IMU/PIR experiment, disabled GPS path, incomplete copied parser.\\
AmberHomeFolder/amber/L-10 and R-11 & Captured Ubuntu runtime evidence for the two AMBER cores; launch configuration binds left/right CAN interfaces, LCM prefixes, UDP ports, solvers, and URDFs. Treat the broader home snapshot as private forensic material, not a deployable image.\\
AmberHomeFolder/amber/sin\_wave/rawLcm & Source definitions for joint targets, arm feedback, Cartesian tasks, plans, solver responses, and gripper control.\\
Amber URDF/amber\_b1.urdf & Independent seven-joint single-arm geometry and model limits; requires measured ROB mounting and tool transforms.\\
Cerebro/Amber-PythonAPI/Amber V2 API & Packed UDP reference client for status, motion, mode, calibration, and gripper commands; requires protocol hardening before production authority.\\
ROBOT Build/Hengdrive*.pdf & Restricted files; contents excluded pending publication clearance.\\
ROBOT Build/B-COMMAND*.pdf & Restricted file; contents excluded and associated installation unproven.\\
ORobotics/media/gallery-originals & Private high-resolution construction and event photographs; strip metadata and review releases.\\
ORobotics/data/gallery\_metadata.json & Visual descriptions useful for search and draft captions, not authoritative part identification.\\
Cerebro/Cerebro/ROBSerialBox.m & Physical serial adapter, base heartbeat, servo/arm helpers.\\
Cerebro/Cerebro/ROBMainViewController.mm & Operator coordination, controller/action/autonomy integration.\\
Cerebro/Cerebro/CameraManager.swift & Camera, RGB-D helper, and bounded frame delivery.\\
Cerebro/Cerebro/ROBAutonomyCoordinator.swift & Session-bounded local autonomy.\\
Cerebro/Cerebro/ROBAI.swift & Optional multimodal Gemini path.\\
Cerebro/Cerebro/GeminiRoboticsProtocol.swift & Versioned high-level action protocol.\\
\bottomrule
\end{tabularx}

\section*{Advanced software source map}

\begin{tabularx}{\textwidth}{>{\ttfamily\scriptsize}p{2.65in} X}
\toprule
RELATIVE PATH & PUBLICATION USE\\
\midrule
Cerebro/Cerebro/ROBSceneSnapshot.swift & Typed people, object, free-space, arm/camera-pose, quality, and confidence boundary for models.\\
Cerebro/Cerebro/ROBMLXRuntime.swift & Actor-owned local LLM, VLM sampling, embeddings, bounded semantic memory, and diagnostics.\\
Cerebro/Cerebro/ROBVideoProtocol.swift & RVID/RBVD video framing, negotiation, limits, codec configuration, and access units.\\
Cerebro/Cerebro/ROBCameraH264Encoder.swift & VideoToolbox NV12 scaling, real-time H.264 encoding, SPS/PPS, AVCC, and recovery output.\\
Cerebro/Cerebro/ROBVideoServer.swift & Authenticated session-bound video subscriptions, bounded send state, feedback, and liveness.\\
Cerebro/Cerebro/ROBStageShow*.swift & Strict show schema, cancellable cue state machine, local/cloud fallback, and diagnostics.\\
Cerebro/Cerebro/ROBSaberChoreography.swift & Allow-listed progressive training transforms and duration profiles.\\
Cerebro/Cerebro/KeyframeAnimationManager.swift & Persisted animation data.\\
Cerebro/docs & Current transport, video, depth, autonomy, and Gemini design notes.\\
ROBController/Consciousness & iPhone/iPad UI, control payload, pairing, Watch relay, action console.\\
ROBController/Consciousness-Watch Watch App & Watch user interface and relay intent.\\
ROBControllerVision & Vision Pro controller, simulator, secure adapter, video decoder, tests.\\
ROBControllerVision/.../ROBVideoPipeline & H.264 validation, SPS/PPS format construction, AVCC sample buffers, renderer recovery, and bounded synthetic fixtures.\\
M2M1-RPLIDAR-iOS-MacOS-Catalyst-/RPLidar & Slamware acquisition and typed lidar publisher.\\
\bottomrule
\end{tabularx}

\chapter{Photo and artwork credits}

All construction photographs in this edition originate from the private Orbitus Robotics gallery archive and are credited to Rodolfo Aramayo / Orbitus Robotics unless the builder supplies a different credit. The frontispiece and software-control illustrations are original generated artwork derived from ROB reference photographs and created for this book set. They contain no official franchise logo, character, or copied page design and are labeled as illustrations rather than engineering evidence.

The book's visual system uses licensed system typefaces and an original weathered retro-space workshop palette. The project is independent and is not affiliated with or endorsed by any film studio, franchise owner, component vendor, or platform company.

Before public printing, complete image releases, confirm third-party product-image and datasheet permissions, remove private metadata and secrets, and obtain a final legal/trademark review appropriate to the distribution plan.

\vfill
\begin{missionbox}[title={NEXT ITERATION}]
Work through the Builder input ledger with the machine physically available. Replace placeholders with measured values, drawings, test reports, verified captions, and contributor stories. Then tag the source snapshot, rebuild all nine PDFs, and perform a page-by-page technical and print proof before the Maker Faire print run.
\end{missionbox}

\ROBSource{This field manual was synthesized from the local ROB presentation, private photo originals and metadata, Arduino firmware and hardware references, and current Cerebro, ROBController, ROBControllerVision, AMBER, and RPLidar repositories. The software consolidation and advanced video pass was completed on 2026-08-11.}

\end{document}
